%قبل از شروع مسئله ترجیحاً چیزی ننویسید. اما اگر خواستید بنویسید، در قالب لیستی (فقرات و شمارش) بنویسید.
%برای اضافه کردن نمره، بعدی از شروع‌مسئله، [x] را اضافه کنید و به جای x  نمره مورد نظر را قرار دهید. برای زیرمسئله‌ها هم به همین صورت است.
%\شروع‌مسئله[10]
\شروع‌مسئله[20] در این سوال باید پیچیدگی الگوریتم‌ها را بررسی کنید. برای هر کدام از الگوریتم‌ها، مشخص کردن مرتبۀ زمانی کفایت می‌کند و نیاز به نوشتن توضیحات نیست. (توجه کنید بررسی پیچیدگی زمانی مستقل از درستی یا غلط بودن آن است.)
\شروع{زیرمسائل}

\زیرمسئله[4] کدام گزینه پیچیدگی زمانی الگوریتم ۱
%\lr{algorithm1}
را در بدترین حالت نشان می‌دهد؟
\begin{large}
	\begin{center}
		\begin{tabular}{c c c}
			$\Theta(n^2)$ & $\Theta(n \log n)$ & $\Theta(\log n)$\\
			$\Theta(2^n)$ &	$\Theta(n\log^2 n)$ & $\Theta(n)$\\
		\end{tabular}
	\end{center}
\end{large}

\پاسخ{
		پیچیدگی زمانی الگوریتم ۱ برابر است با 
	$\Theta(n\log n)$.
	
}
%\شروع{شمارش}
%\فقره
%$\Theta(n^2)$
%
%\فقره
%$\Theta(n \log n)$
%
%\فقره
%$\Theta(\log n)$
%
%\فقره
%$\Theta(2^n)$
%
%\فقره
%$\Theta(n\log^2 n)$
%
%\فقره
%$\Theta(n)$
%
%\پایان{شمارش}


\زیرمسئله[4] کدام گزینه پیچیدگی زمانی الگوریتم ۲
%\lr{algorithm2}
را در بدترین حالت نشان می‌دهد؟

\begin{large}
	\begin{center}
		\begin{tabular}{c c c}
			$\Theta(n^2)$ & $\Theta(n \log n)$ & $\Theta(\log n)$\\
			$\Theta(2^n)$ &	$\Theta(n\log^2 n)$ & $\Theta(n)$\\
		\end{tabular}
	\end{center}
\end{large}

\پاسخ{
		پیچیدگی زمانی الگوریتم ۲ برابر است با 
	$\Theta(n^2)$.
	
}

%\شروع{شمارش}
%\فقره
%$\Theta(n^2)$
%
%\فقره
%$\Theta(n \log n)$
%
%\فقره
%$\Theta(\log n)$
%
%\فقره
%$\Theta(2^n)$
%
%\فقره
%$\Theta(n\log^2 n)$
%
%\فقره
%$\Theta(n)$
%
%
%\پایان{شمارش}


\زیرمسئله[4] کدام گزینه پیچیدگی زمانی الگوریتم ۳
%\lr{algorithm3}
را در بدترین حالت نشان می‌دهد؟
\begin{large}
	\begin{center}
		\begin{tabular}{c c c}
			$\Theta(n^2)$ & $\Theta(n \log n)$ & $\Theta(\log n)$\\
			$\Theta(2^n)$ &	$\Theta(n\log^2 n)$ & $\Theta(n)$\\
		\end{tabular}
	\end{center}
\end{large}

\پاسخ{
		پیچیدگی زمانی الگوریتم ۳ برابر است با 
	$\Theta(n)$.
	فرض کنید تعداد سطرها برابر $m$ و تعداد ستون‌ها برابر ‌$n$ باشد. در این صورت با توجه به خط ۸۲ و ۸۴ در کد
	که نشان دهندۀ این موضوع است که هزینۀ انجام این کار برابر است با 
	$\Theta(n+m)$،
	بنابراین رابطۀ زیر برقرار است:
	$$T(m,n ) = \Theta(m+n) +T(m/2, n/2)$$
	به طور بازگشتی، از خود رابطه اصلی در خودش استفاده می‌کنیم، به این ترتیب داریم:
	$$ T(m,n ) = c(m+n) + c(\frac{m+n}{2}) + c(\frac{m+n}{4}) + c(\frac{m+n}{8}) + ... $$
	بنابراین از طرفی می‌دانیم که $T(m,n)$ از 
	$2c(m+n)$
	کم‌تر است و هم‌چنین از ‌
	$c(m+n)$
	بیشتر است. بنابراین داریم: 
	$$T(m,n) = \Theta(m+n)$$
	و به طور خاص داریم 
	$T(n,n) = \Theta(n)$.
	
}
%\شروع{شمارش}
%\فقره
%$\Theta(n^2)$
%
%\فقره
%$\Theta(n \log n)$
%
%\فقره
%$\Theta(\log n)$
%
%\فقره
%$\Theta(2^n)$
%
%\فقره
%$\Theta(n\log^2 n)$
%
%\فقره
%$\Theta(n)$
%
%
%\پایان{شمارش}

\زیرمسئله[4] کدام گزینه پیچیدگی زمانی الگوریتم ۴
%\lr{algorithm4}
را در بدترین حالت نشان می‌دهد؟

\begin{large}
	\begin{center}
		\begin{tabular}{c c c}
			$\Theta(n^2)$ & $\Theta(n \log n)$ & $\Theta(\log n)$\\
			$\Theta(2^n)$ &	$\Theta(n\log^2 n)$ & $\Theta(n)$\\
		\end{tabular}
	\end{center}
\end{large}

\پاسخ{
		پیچیدگی زمانی الگوریتم ۴ برابر است با 
	$\Theta(n)$.
	
}

%\شروع{شمارش}
%\فقره
%$\Theta(n^2)$
%
%\فقره
%$\Theta(n \log n)$
%
%\فقره
%$\Theta(\log n)$
%
%\فقره
%$\Theta(2^n)$
%
%\فقره
%$\Theta(n\log^2 n)$
%
%\فقره
%$\Theta(n)$
%
%
%\پایان{شمارش}


\زیرمسئله[4] کدام گزینه پیچیدگی زمانی الگوریتم ۵
%\lr{algorithm5}
را در بدترین حالت نشان می‌دهد؟
\begin{large}
	\begin{center}
		\begin{tabular}{c c c}
			$\Theta(n^2)$ & $\Theta(n \log n)$ & $\Theta(\log n)$\\
			$\Theta(2^n)$ &	$\Theta(n\log^2 n)$ & $\Theta(n)$\\
		\end{tabular}
	\end{center}
\end{large}

\پاسخ{
		پیچیدگی زمانی الگوریتم ۵ برابر است با 
	$\Theta(n)$.
	ابعاد ماتریس اولیه را $n\times n$
	می‌گیریم. با این فرض، حلقۀ اول و دوم هر کدام 
	$n$
	بار تکرار می‌شوند. همچنین پیدا کردن عنصر بیشینه نیز با تعداد خانه‌ها متناسب است که از 
	$6n$
	کمتر است. همچنین غیر از روند بازگشتی الگوریتم، بقیۀ اعمال در زمان ثابت انجام می‌گیرند. پس برای عدد ثابت  به اندازۀ کافی بزرگ $c$ داریم:
	\begin{equation*}
	T(n) \leq cn + T(\frac{n - 3}{2})
	\end{equation*}
	که نشان می‌دهد الگوریتم از مرتبه $n$ است (می‌توان از استقرا برای اثبات این موضوع استفاده کرد).

}
%\شروع{شمارش}
%\فقره
%$\Theta(n^2)$
%
%\فقره
%$\Theta(n \log n)$
%
%\فقره
%$\Theta(\log n)$
%
%\فقره
%$\Theta(2^n)$
%
%\فقره
%$\Theta(n\log^2 n)$
%
%\فقره
%$\Theta(n)$
%
%
%\پایان{شمارش}

\پایان{زیرمسائل}
%اگر مسئله چند قسمت داشت به این صورت عمل کنید:
%\شروع{زیرمسائل}
%\زیرمسئله[0] % مسئلۀ ۱آ
%\زیرمسئله % مسئلۀ ۱ب
%\زیرمسئله % مسئلۀ ۱پ
%\زیرمسئله % مسئلۀ ۱ت
%\پایان{زیرمسائل}





%%%%%%%%%%%%%%%%%%%%%%%%%%%%%%%%%%%%%%%%%%%%%%%%%%%%%
% برخی دستورات مفید برای تمرین‌ها. %
%%%%%%%%%%%%%%%%%%%%%%%%%%%%%%%%%%%%%%%%%%%%%%%%%%%%%

% فرمول‌های نمونه:
%$f(x) = \Theta(x)$
%$T(x, y) \leq \log(x) + 2^y + \binom{2n}{n}$
% {\tt code\_function}


% لیست موارد بدون شماره‌گذاری:
%\شروع{فقرات}
%    \فقره مورد اول 
%        \شروع{فقرات}
%            \فقره مورد اول داخلی 
%                \شروع{فقرات}
%                    \فقره اولین مورد خیلی داخلی 
%                    \فقره دومین مورد خیلی داخلی 
%                \پایان{فقرات}
%            \فقره مورد دوم داخلی 
%        \پایان{فقرات}
%    \فقره مورد دوم 
%\پایان{فقرات}

% لیست موارد شماره‌گذاری شده:
%\شروع{شمارش}
%    \فقره یک مورد 
%    \فقره یک مورد دیگر 
%    \فقره باز هم یک مورد دیگر
%\پایان{شمارش}

% برای نوشتن یک فرمول چند خطی که یک شماره بخوره. برای حالت‌های دیگر به سایت زیر مراجعه کنید:
% https://www.overleaf.com/learn/latex/Aligning_equations_with_amsmath
%\begin{align} 
%    \begin{split}
%        (x+y)^3 &= (x+y)^2(x+y) \\
%                &= (x^2+2xy+y^2)(x+y) \\
%                &= (x^3+2x^2y+xy^2) + (x^2y+2xy^2+y^3) \\
%                &= x^3+3x^2y+3xy^2+y^3
%    \end{split}                                 
%\end{align}

% نوشتن ماتریس و آرایه دوبعدی
%\begin{align}
%    A = 
%    \begin{bmatrix}
%        A_{11} & A_{21} \\
%        A_{21} & A_{22}
%    \end{bmatrix}
%\end{align}

% اضافه کردن تصویر با شرح. برای اضافه کردن تصویر بدون شرح، فقط دستور درج‌تصویر را استفاده کنید.
%\شروع{شکل}[h!]
%\تنظیم‌ازوسط
%  \درج‌تصویر[پهنا=0.5\پهنای‌متن]{img.jpg}
%\شرح{شرح تصویر}
%\برچسب{تصویر:۱}
%\پایان{شکل}
%\پایان{زیرمسائل}

%کد زیر را ببینید.
%%\lstinputlisting{source_filename.py}
%\begin{latin}
%\begin{lstlisting}
%# Import the modules
%import sys
%import random
%
%ans = True
%
%while ans:
%question = raw_input("Ask the magic 8 ball a question: (press enter to quit) ")
%
%answers = random.randint(1,8)
%
%if question == "":
%sys.exit()
%
%elif answers == 1:
%print "It is certain"
%
%elif answers == 2:
%print "Outlook good"
%
%elif answers == 3:
%print "You may rely on it"
%
%elif answers == 4:
%print "Ask again later"
%
%elif answers == 5:
%print "Concentrate and ask again"
%
%elif answers == 6:
%print "Reply hazy, try again"
%
%elif answers == 7:
%print "My reply is no"
%
%elif answers == 8:
%print "My sources say no"
%\end{lstlisting}
%\end{latin}
